% !TEX program = xelatex
% =====================================================================
% Channel Proxy v4 -> v6 -> v7 修订记录
% IMPORTANT: 必须用 XeLaTeX 编译 (Overleaf: Menu -> Compiler -> XeLaTeX)
% =====================================================================
\documentclass[11pt,a4paper,UTF8,fontset=fandol]{ctexart}
\usepackage[margin=2.2cm]{geometry}
\usepackage{listings}
\usepackage[table]{xcolor}
\usepackage{booktabs}
\usepackage{longtable}
\usepackage{array}
\usepackage{amsmath}
\usepackage{amssymb}
\usepackage{hyperref}
\usepackage{enumitem}
\usepackage{titlesec}
\usepackage{caption}

\hypersetup{
  colorlinks=true,
  linkcolor=black!50!blue,
  urlcolor=blue!60!black,
  citecolor=black,
}

% --------------- code style ---------------
\definecolor{codegreen}{rgb}{0.20,0.55,0.20}
\definecolor{codeblue}{rgb}{0.10,0.20,0.65}
\definecolor{codered}{rgb}{0.75,0.10,0.10}
\definecolor{codebg}{rgb}{0.975,0.975,0.975}
\definecolor{starcolor}{rgb}{0.92,0.50,0.05}
\definecolor{rowblue}{rgb}{0.92,0.95,1.00}
\definecolor{rowgray}{rgb}{0.96,0.96,0.96}

\lstdefinestyle{pystyle}{
    backgroundcolor=\color{codebg},
    commentstyle=\color{codegreen}\itshape,
    keywordstyle=\color{codeblue}\bfseries,
    stringstyle=\color{codered},
    basicstyle=\ttfamily\footnotesize,
    breakatwhitespace=false,
    breaklines=true,
    captionpos=b,
    keepspaces=true,
    numbers=left,
    numberstyle=\tiny\color{black!50},
    numbersep=6pt,
    showspaces=false,
    showstringspaces=false,
    showtabs=false,
    tabsize=2,
    language=Python,
    frame=single,
    framerule=0.4pt,
    rulecolor=\color{black!30},
    aboveskip=6pt,
    belowskip=6pt,
}
\lstset{style=pystyle}

\renewcommand{\lstlistingname}{代码}

% --------------- star severity ---------------
\newcommand{\sone}{\textcolor{starcolor}{$\bigstar$}}
\newcommand{\stwo}{\textcolor{starcolor}{$\bigstar\bigstar$}}
\newcommand{\sthree}{\textcolor{starcolor}{$\bigstar\bigstar\bigstar$}}

\newcommand{\vfour}{\texttt{v4.py}}
\newcommand{\vsix}{\texttt{v6.py}}
\newcommand{\vseven}{\texttt{v7.py}}

% --------------- bug header ---------------
\newcommand{\bugitem}[3]{%
  \subsection{[\##1] #2 \hfill #3}%
  \vspace{-0.5em}%
}

% --------------- title ---------------
\titleformat{\section}
  {\normalfont\Large\bfseries}{\thesection.}{0.6em}{}
\titleformat{\subsection}
  {\normalfont\large\bfseries}{\thesubsection}{0.6em}{}

\setlist[itemize]{leftmargin=1.4em,topsep=2pt,itemsep=2pt}

\title{\textbf{G1C Multi-UE MIMO Channel Proxy} \\[0.4em]
       \large 修订记录: v4 $\to$ v6 $\to$ v7}
\author{}
\date{2026-05-08}

% =====================================================================
\begin{document}
\maketitle

\begin{abstract}
\noindent
本文档汇总 \vfour{} $\to$ \vsix{} $\to$ \vseven{} 的全部代码修订, 共 29 条.
v6 集中处理 v4 中的信道生成物理建模与系统鲁棒性问题; v7 在 v6 基础上针对 OAI
attach 阶段引入控制面稳定机制并修正 P1B 距离推导. 每条改动按影响范围标注严重
程度: \sthree{} 影响信道输出物理正确性或多 UE 场景下的并发吞吐;
\stwo{} 影响系统鲁棒性, 资源占用或控制面稳定性; \sone{} 属于代码风格,
兼容性或边界情况. 每条改动给出 \vfour{} 与 \vseven{} 中对应位置的代码片段以便核验.
\end{abstract}

\tableofcontents
\vspace{1em}

% =====================================================================
\section{星级说明与版本演进}
\label{sec:overview}

\subsection{星级定义}
\begin{itemize}
  \item \sthree{} \ \textbf{核心修订}: 涉及信道生成的物理建模数学错误,
        或在多 UE 场景下决定整体吞吐的并发瓶颈. 该类问题不修复时, 下游性能
        指标 (NMSE / RSRP / SRS bin) 不应作为基线参考.
  \item \stwo{} \ \textbf{重要修订}: 影响系统鲁棒性, 显存占用,
        异常恢复, 或 OAI 控制面收敛过程. 不修复会限制实验时长或重现性.
  \item \sone{} \ \textbf{常规修订}: 边界检查, 兼容性, 死代码清理,
        日志格式. 不影响实验结论.
\end{itemize}

\subsection{版本演进示意}
\begin{center}
\renewcommand{\arraystretch}{1.25}
\begin{tabular}{>{\bfseries}l p{12cm}}
\toprule
v4.py & 通用通道生成 + IPC ring buffer 框架 \\
\midrule
v6.py & 物理建模 \& 鲁棒性 \& 性能 \& CUDA Graph 修订 (本文 \#1$\sim$\#26) \\
\midrule
v7.py & v6 + 控制面稳定窗口 + 自动 handoff + P1B 距离模式 (本文 \#27$\sim$\#29) \\
\bottomrule
\end{tabular}
\end{center}

\subsection{涉及文件}
\begin{itemize}
  \item \texttt{DevChannelProxyJIN/vRAN\_Socket/G1C\_MultiUE\_MIMO\_Channel\_Proxy/v4.py}
        \quad (3164 行)
  \item \texttt{DevChannelProxyJIN/vRAN\_Socket/G1C\_MultiUE\_MIMO\_Channel\_Proxy/v6.py}
        \quad (3554 行)
  \item \texttt{DevChannelProxyJIN/vRAN\_Socket/G1C\_MultiUE\_MIMO\_Channel\_Proxy/v7.py}
        \quad (3660 行)
  \item \texttt{verify\_issue\_2.py}: 5 阶段物理一致性验证脚本 (附录 \ref{sec:verify})
\end{itemize}

% =====================================================================
\section{修订总览表}
\label{sec:summary}

\begin{center}
\renewcommand{\arraystretch}{1.18}
\rowcolors{2}{rowgray}{white}
\begin{longtable}{c c p{8.6cm} c}
\toprule
\textbf{编号} & \textbf{类别} & \textbf{修订内容} & \textbf{严重度} \\
\midrule
\endfirsthead
\toprule
\textbf{编号} & \textbf{类别} & \textbf{修订内容} & \textbf{严重度} \\
\midrule
\endhead
\bottomrule
\endfoot

1  & 物理 & 载频单位由 \texttt{3.5} 修正为 \texttt{3.5e9} Hz                 & \sthree \\
2  & 物理 & \texttt{sample\_times} 跨 batch 累积递增, 避免输出重复 realization & \sthree \\
3  & 物理 & \texttt{velocities} 改为水平面随机方向 + 标量速度模长 ($V_z=0$)    & \stwo   \\
4  & 物理 & LOS 角度 (aoa/aod/zoa/zod) 由全零改为 3GPP TR38.901 范围内随机    & \stwo   \\
5  & 物理 & \texttt{distance\_3d} 由硬编码 1\,m 改为可配置                     & \stwo   \\
6  & 物理 & 能量归一化由 per-batch 改为 batch-0 全局基准                       & \sthree \\
7  & 物理 & 新增 \texttt{\_validate\_physics\_params()} 在 CLI 解析后强校验    & \sone   \\
8  & 鲁棒 & \texttt{batch\_idx} 在 \texttt{generate\_fn()} 成功后立即递增      & \sone   \\
9  & 鲁棒 & 主循环增加 producer 存活检查, 死亡时 clean shutdown                & \stwo   \\
10 & 鲁棒 & 异常重试加 20 次连续失败上限                                       & \stwo   \\
11 & 鲁棒 & Socket 模式拒绝 \texttt{num\_ues>1} 而非静默广播                   & \stwo   \\
12 & 鲁棒 & \texttt{\_stalled\_ue\_logged} 在 UE 恢复时清除                    & \sone   \\
13 & 鲁棒 & 启动时校验 \texttt{gnb\_ant == gnb\_nx * gnb\_ny}                  & \sone   \\
14 & 鲁棒 & \texttt{random.sample} 越界给出明确错误信息                        & \sone   \\
15 & 性能 & UL superposition 路径 \texttt{skip\_quant=True} 跳过冗余量化       & \stwo   \\
16 & 性能 & \texttt{NoiseProducer} dtype/buffer 缩减 ($\sim 16\times$ VRAM)    & \stwo   \\
17 & 性能 & \texttt{GTBatchSaver} 异步写盘                                     & \stwo   \\
18 & 性能 & \texttt{bypass\_copy} 预分配 + 缓存复用                            & \sone   \\
19 & 性能 & UL \texttt{set\_last\_ul\_rx\_ts} 每 slot 仅 wake 一次             & \sone   \\
20 & 性能 & \texttt{IPCRingBuffer} 同步操作移出锁外                            & \sthree \\
21 & CUDA Graph & 改用 \texttt{streamCaptureModeRelaxed}                       & \sone   \\
22 & CUDA Graph & capture 失败可恢复 (5 次重试) 而非永久禁用                   & \sone   \\
23 & 杂项 & PEP 604 \texttt{Tuple|None} 改为 \texttt{Optional[Tuple[...]]}    & \sone   \\
24 & 杂项 & CPU 模式下 \texttt{noise\_dBFS} 不再静默丢失                       & \sone   \\
25 & 杂项 & 移除未使用的 \texttt{\_ul\_clip\_3d}                               & \sone   \\
26 & 杂项 & CUDA Graph 中 \texttt{pl\_linear} 去除 Python 分支                 & \sone   \\
\midrule
27 & 控制面 & 新增 attach-stable 时间窗冻结 \texttt{sample\_times}              & \stwo   \\
28 & 控制面 & 新增基于触发文件的 auto handoff 机制                              & \sone   \\
29 & 物理 & P1B 距离改为可选模式, 默认固定值代替 $\tau$ 推导                   & \stwo   \\
\end{longtable}
\end{center}

% =====================================================================
\section{物理正确性}
\label{sec:physics}

% ---------------------------------------------------------------------
\bugitem{1}{载频单位 \texttt{3.5} $\to$ \texttt{3.5e9}}{\sthree}

\textbf{现象.} \vfour{} 中载频常量 \texttt{carrier\_frequency} 写为标量
\texttt{3.5}, 单位语义缺失. Sionna 的 \texttt{PanelArray} 与
\texttt{ChannelCoefficientsGeneratorJIN} 期望该参数以 Hz 为单位,
被解释为 \texttt{3.5\,Hz} 时:
\begin{itemize}
  \item 波长 $\lambda_0 = c/f = 8.57\times 10^{7}\,\text{m}$
        (期望约 \,$0.0857\,\text{m}$);
  \item Doppler 项 $\exp(j\,2\pi v \hat{r}\cdot t / \lambda)$ 在合理仿真时长内
        相位旋转量趋于零, 信道时变性消失;
  \item 天线间距按 $\lambda/2$ 计算时, 物理尺度差 9 个数量级.
\end{itemize}

\textbf{\vfour{} 原始代码.}
\begin{lstlisting}[caption={v4.py: line 194}]
# OFDM NR numerology=1
carrier_frequency = 3.5
\end{lstlisting}

\textbf{\vseven{} 修订.}
\begin{lstlisting}[caption={v7.py: line 234}]
# OFDM NR numerology=1
carrier_frequency = 3.5e9
\end{lstlisting}

\textbf{影响范围.} 该项与 \#2 协同导致 \vfour{} producer 输出的信道实质上为
静态 realization. 不修复时, 下游 NMSE 与 RSRP 测量不具有物理意义.

% ---------------------------------------------------------------------
\bugitem{2}{\texttt{sample\_times} 跨 batch 累积递增}{\sthree}

\textbf{现象.} \texttt{\_H\_TTI\_sequential\_fft\_o\_ELW2\_noProfile()}
是无内部随机源的纯确定性函数, 输出仅由 (topology, sample\_times, h\_field, aoa, zoa)
决定. \vfour{} 在循环外构造一次 \texttt{sample\_times\_fixed} 后, 主循环每次以
相同入参调用, 因此 producer 在整个生命周期内输出完全相同的一份 channel realization.

\textbf{\vfour{} 原始代码.}
\begin{lstlisting}[caption={v4.py: lines 1817--1819 (定义于循环外, 循环内不变)}]
sample_times_fixed = _tf.cast(
    _tf.range(buffer_symbol_size), gen.rdtype
) / _tf.constant(params['scs'], gen.rdtype)
\end{lstlisting}

\textbf{\vseven{} 修订.} 使用 \texttt{int64} 基底跨 batch 递增,
避免 batch 计数过大后浮点累加丢精度.
\begin{lstlisting}[caption={v7.py: lines 2088--2092}]
def _build_sample_times(b_idx):
    # absolute symbol indices: [b_idx*B, b_idx*B+1, ..., b_idx*B+B-1]
    start = _tf.constant(b_idx * _bss, dtype=_tf.int64)
    idx_i64 = start + _tf.range(_bss, dtype=_tf.int64)
    return _tf.cast(idx_i64, gen.rdtype) / _scs_tf
\end{lstlisting}

\textbf{与 \#1 的耦合.} \#1 与 \#2 互为掩盖: 仅修复 \#2 时由于
$\lambda \sim 10^{7}\,\text{m}$, Doppler 仍近似为 0; 仅修复 \#1 时由于
sample\_times 不变, 输出仍重复. 必须同时修复后,
跨 batch 信道演化才会显现 (相对差异约 93.5\%, 见附录 \ref{sec:verify}).

% ---------------------------------------------------------------------
\bugitem{3}{\texttt{velocities} 物理建模}{\stwo}

\textbf{现象.} \vfour{} 对 3 个速度分量分别独立采样
$|\mathcal{N}(\text{Speed}, 0.1)|$, 总速度模长
$|\boldsymbol{v}| \approx \sqrt{3}\cdot\text{Speed}$ 而非
$\text{Speed}$ 本身; 同时 \texttt{tf.abs} 将所有分量限制于第一象限,
方向分布严重失真.

\textbf{\vfour{} 原始代码.}
\begin{lstlisting}[caption={v4.py: lines 1774--1775}]
velocities = _tf.abs(_tf.random.normal(
    shape=[batch_size, N_UE, 3], mean=Speed_local, stddev=0.1, dtype=_tf.float32))
\end{lstlisting}

\textbf{\vseven{} 修订.} 水平面方向角 $\phi \sim \mathcal{U}[0, 2\pi)$,
速度模长 $|v| = |\mathcal{N}(\text{Speed}, 0.1)|$, 垂直分量 $V_z = 0$.
\begin{lstlisting}[caption={v7.py: lines 2011--2023}]
_phi = _tf.random.uniform(
    shape=[batch_size, N_UE], minval=0.0, maxval=2 * float(PI),
    dtype=_tf.float32)
_speed_mag = _tf.abs(_tf.random.normal(
    shape=[batch_size, N_UE], mean=Speed_local, stddev=0.1,
    dtype=_tf.float32))
velocities = _tf.stack([
    _speed_mag * _tf.cos(_phi),                       # Vx
    _speed_mag * _tf.sin(_phi),                       # Vy
    _tf.zeros_like(_speed_mag),                       # Vz=0 (planar)
], axis=-1)
\end{lstlisting}

% ---------------------------------------------------------------------
\bugitem{4}{LOS 路径角度初始化}{\stwo}

\textbf{现象.} \vfour{} 将所有 BS-UE 对的 LOS azimuth/zenith
(\texttt{los\_aoa, los\_aod, los\_zoa, los\_zod}) 初始化为零,
所有 UE 共享同一 LOS 方向, 失去空间多样性.

\textbf{\vfour{} 原始代码.}
\begin{lstlisting}[caption={v4.py: lines 1776--1779}]
los_aoa = _tf.zeros([batch_size, N_BS, N_UE])
los_aod = _tf.zeros([batch_size, N_BS, N_UE])
los_zoa = _tf.zeros([batch_size, N_BS, N_UE])
los_zod = _tf.zeros([batch_size, N_BS, N_UE])
\end{lstlisting}

\textbf{\vseven{} 修订.} 按 3GPP TR38.901 范围采样:
azimuth $\in [0, 2\pi)$, zenith $\in [\pi/3, 2\pi/3]$ (水平面 $\pm 30^{\circ}$).
\begin{lstlisting}[caption={v7.py: lines 2030--2041}]
los_aoa = _tf.random.uniform(
    shape=[batch_size, N_BS, N_UE], minval=0.0, maxval=2 * float(PI),
    dtype=_tf.float32)
los_aod = _tf.random.uniform(
    shape=[batch_size, N_BS, N_UE], minval=0.0, maxval=2 * float(PI),
    dtype=_tf.float32)
los_zoa = _tf.random.uniform(
    shape=[batch_size, N_BS, N_UE],
    minval=float(PI) / 3, maxval=2 * float(PI) / 3, dtype=_tf.float32)
los_zod = _tf.random.uniform(
    shape=[batch_size, N_BS, N_UE],
    minval=float(PI) / 3, maxval=2 * float(PI) / 3, dtype=_tf.float32)
\end{lstlisting}

% ---------------------------------------------------------------------
\bugitem{5}{\texttt{distance\_3d} 硬编码 1\,m}{\stwo}

\textbf{现象.} \vfour{} 将所有 BS-UE 距离硬编码为 1\,m,
脱离任何合理蜂窝几何, 路径损耗模型失真.

\textbf{\vfour{} 原始代码.}
\begin{lstlisting}[caption={v4.py: line 1782}]
distance_3d = _tf.ones([1, N_BS, N_UE])
\end{lstlisting}

\textbf{\vsix{} 修订}: 当 P1B 数据可用时由最短路径 $\tau$ 推算 $d \approx \tau \cdot c$
(下限 1\,m), 否则使用 \texttt{--default-distance-m} 配置值. 该方案在实际 P1B 数据上
表现不佳, 详见 \#29.

\textbf{\vseven{} 修订}: 见 \#29 (默认改为固定 100\,m, $\tau$ 推导改为可选模式).

% ---------------------------------------------------------------------
\bugitem{6}{Energy normalization 抹去 Rayleigh 功率波动}{\sthree}

\textbf{现象.} \vfour{} 对每个 batch 各自计算总能量并归一化,
等效于强制每 batch 总能量恒等. Rayleigh fading 的功率起伏被压平,
Doppler 引入的小尺度功率波动同时被消去, 下游 NMSE 与 RSRP 的统计意义失效.

\textbf{\vfour{} 原始代码.}
\begin{lstlisting}[caption={v4.py: lines 1834--1835}]
energy = _tf.reduce_sum(_tf.abs(h_c128) ** 2, axis=3, keepdims=True)
h_norm = h_delay / _tf.cast(_tf.sqrt(energy + 1e-30), h_delay.dtype)
\end{lstlisting}

\textbf{\vseven{} 修订.} 在 batch 0 计算一次全局标量 $|H|_{\text{rms}}$,
所有后续 batch 共除该常数.
\begin{lstlisting}[caption={v7.py: lines 2117--2125}]
_h0 = _generate_unnorm(_build_sample_times(0))
_ref_energy = _tf.reduce_mean(_tf.abs(_h0) ** 2, keepdims=False)
_ref_norm = _tf.cast(_tf.sqrt(_ref_energy + 1e-30), _h0.dtype)
print(f"[v7 ChannelProducer] reference energy normalization: "
      f"|H|_rms = {float(_tf.abs(_ref_norm)):.4e}")

def _generate_eager(sample_times):
    h_c128 = _generate_unnorm(sample_times)
    return h_c128 / _ref_norm   # global, time-invariant scaling
\end{lstlisting}

副作用: 由于不再做 per-batch 归一化, 每个 batch 之间能量自然连续,
Phase C 验证显示批边界跳变约 $2.6\times 10^{-7}$ (浮点精度级).

% ---------------------------------------------------------------------
\bugitem{7}{物理参数边界校验}{\sone}

\textbf{现象.} \vfour{} 不校验 \texttt{carrier\_frequency}, \texttt{scs},
\texttt{Speed}; 用户传入异常值时静默执行.

\textbf{\vseven{} 修订.} 新增 \texttt{\_validate\_physics\_params()},
在 CLI 解析后调用, 使用 \texttt{ValueError} 而非 \texttt{assert}
(\texttt{python -O} 下 assert 会被剥除).
\begin{lstlisting}[caption={v7.py: lines 256--274}]
def _validate_physics_params(cf=None, sc=None, sp=None):
    """Validate physics parameters. Uses module globals if args omitted.
    Raises ValueError (not assert) so it survives `python -O`."""
    cf = carrier_frequency if cf is None else cf
    sc = scs if sc is None else sc
    sp = Speed if sp is None else sp
    _wavelength = 3e8 / cf
    if not (1e8 < cf < 1e12):
        raise ValueError(
            f"carrier_frequency={cf} Hz - "
            f"expected 100 MHz - 1 THz for 5G/sub-THz")
    if not (1e-4 < _wavelength < 10):
        raise ValueError(
            f"wavelength={_wavelength} m - "
            f"implausible for 5G (expected mm-cm range)")
    if sc not in (15e3, 30e3, 60e3, 120e3, 240e3):
        raise ValueError(
            f"scs={sc} Hz - not a standard NR numerology")
\end{lstlisting}

% =====================================================================
\section{鲁棒性}
\label{sec:robustness}

% ---------------------------------------------------------------------
\bugitem{8}{\texttt{batch\_idx} 递增时机}{\sone}

\textbf{现象.} \vfour{} 将计数 (\texttt{symbol\_counter}) 放在循环末尾递增.
若 downstream (FFT, ring buffer 写入) 抛异常, 重试时计数语义混乱.

\textbf{\vseven{} 修订.} 在 \texttt{generate\_fn()} 成功返回后立即递增
\texttt{batch\_idx}, 与下游处理解耦.
\begin{lstlisting}[caption={v7.py: lines 2199--2202}]
# v7 #V6-4: increment batch_idx IMMEDIATELY after successful
# generate_fn() so a downstream exception does not trigger a
# duplicate batch on retry.
batch_idx += 1
\end{lstlisting}

% ---------------------------------------------------------------------
\bugitem{9}{Producer 存活检查}{\stwo}

\textbf{现象.} \vfour{} 仅在 warmup 阶段检查 producer 进程存活,
主循环进入后不再检查. Producer 异常退出后, 主进程持续从 ring buffer 读取
stale data, 错误以静默形式向下游传播.

\textbf{\vsix/\vseven{} 修订.} 主循环中每 1000 次迭代检查 \texttt{proc.is\_alive()},
死亡时打印 exit code 并 break. 同时 ring buffer 读取增加超时保护.

% ---------------------------------------------------------------------
\bugitem{10}{异常重试上限}{\stwo}

\textbf{现象.} \vfour{} 在 producer 主循环 \texttt{except Exception} 后
\texttt{traceback.print\_exc()} 并 \texttt{time.sleep(1.0)}, 不设上限.
持续故障会以每秒一份 traceback 的频率污染日志, 同时占用 GPU 资源.

\textbf{\vfour{} 原始代码.}
\begin{lstlisting}[caption={v4.py: lines 1880--1884}]
except Exception as e:
    print(f"[v4 UnifiedChannelProducerProcess] ERROR in generation loop: {e}")
    import traceback
    traceback.print_exc()
    time.sleep(1.0)
\end{lstlisting}

\textbf{\vseven{} 修订.}
\begin{lstlisting}[caption={v7.py: lines 2154--2164}]
except Exception as e:
    consecutive_errors += 1
    print(f"[v7 UnifiedChannelProducerProcess] ERROR "
          f"({consecutive_errors}/{MAX_CONSECUTIVE_ERRORS}): {e}")
    import traceback
    traceback.print_exc()
    if consecutive_errors >= MAX_CONSECUTIVE_ERRORS:
        print(f"[v7 UnifiedChannelProducerProcess] Too many "
              f"consecutive errors - exiting.")
        break
    time.sleep(1.0)
\end{lstlisting}
成功一次后 \texttt{consecutive\_errors} 重置为 0, 因此偶发故障不会累计.

% ---------------------------------------------------------------------
\bugitem{11}{Socket 模式 multi-UE 处理}{\stwo}

\textbf{现象.} \vfour{} socket 模式下若 \texttt{num\_ues>1},
将同一份 processed signal 广播给所有 UE; 结果错误但无任何提示.

\textbf{\vseven{} 修订.} 启动时显式拒绝:
\begin{verbatim}
if num_ues > 1 and mode == "socket":
    raise ValueError("Socket mode does not support num_ues>1")
\end{verbatim}

% ---------------------------------------------------------------------
\bugitem{12}{\texttt{\_stalled\_ue\_logged} 状态恢复}{\sone}

\textbf{现象.} \vfour{} UE 一旦被标记 stalled, 永远不会移出
\texttt{\_stalled\_ue\_logged}; UE 实际恢复后既无日志也无后续检测.

\textbf{\vseven{} 修订.}
\begin{verbatim}
_recovered = self._stalled_ue_logged & active_set
if _recovered:
    print(f"UE recovered from stall: {sorted(_recovered)}")
    self._stalled_ue_logged -= _recovered
\end{verbatim}

% ---------------------------------------------------------------------
\bugitem{13}{天线数一致性校验}{\sone}

\textbf{现象.} 用户传入 \texttt{--gnb-ant 4 --gnb-nx 1 --gnb-ny 1} 时,
\texttt{PanelArray} 维度推导失败, 报错信息晦涩.

\textbf{\vseven{} 修订.} \texttt{main()} 入口处显式校验
\texttt{gnb\_ant == gnb\_nx * gnb\_ny}, 不通过则直接报错并指明具体值.

% ---------------------------------------------------------------------
\bugitem{14}{\texttt{random.sample} 边界检查}{\sone}

\textbf{现象.} \texttt{num\_ues > len(all\_indices)} 时,
\vfour{} 直接命中 stdlib \texttt{ValueError}, 用户难以定位是 P1B 索引耗尽.

\textbf{\vseven{} 修订.} 提前检查并给出可定位错误:
\begin{verbatim}
if num_ues > len(all_indices):
    raise ValueError(
        f"--num-ues={num_ues} exceeds available RX indices "
        f"({len(all_indices)}) in {p1b_npz}")
\end{verbatim}

% =====================================================================
\section{性能与内存}
\label{sec:performance}

% ---------------------------------------------------------------------
\bugitem{15}{UL 路径 \texttt{skip\_quant}}{\stwo}

\textbf{现象.} UL superposition 中各 UE 的 \texttt{process\_slot\_ipc()}
调用结束时, 实际向上累加用的是复数 \texttt{gpu\_out}, 但 \vfour{} 仍执行了
完整的 clip + cast$\to$int16 写入 \texttt{\_ul\_dummy\_out};
该输出立即被丢弃, 浪费一次 GPU memcpy + kernel.

\textbf{\vseven{} 修订.} 新增 \texttt{skip\_quant} 参数, UL 路径调用时
设置 \texttt{skip\_quant=True}; 对应 CUDA Graph 也跳过最终量化分支
(\texttt{\_use\_graph\_this\_call = not skip\_quant}).

% ---------------------------------------------------------------------
\bugitem{16}{\texttt{NoiseProducer} 显存占用}{\stwo}

\textbf{现象.}
\begin{center}
\begin{tabular}{l c c c}
\toprule
版本 & dtype & BATCH\_SIZE & maxlen \\
\midrule
v4 & float64 & 64 & 256 \\
v7 & float32 & 32 & 64 \\
\bottomrule
\end{tabular}
\end{center}
v4 配置下, 单个 noise buffer 占约 480\,MB (\texttt{gnb\_ant=4}),
4 UE $\times$ 2 (DL+UL) 接近 4\,GB.

\textbf{\vseven{} 修订.} dtype 降为 float32, BATCH\_SIZE 降为 32.
单 buffer 约 15\,MB, 4 UE $\times$ 2 约 120\,MB ($\sim 16\times$ VRAM 节省).
噪声物理统计仅需 $\sim 5$ bit 精度, 远小于 float64 的 52 bit; consumer 端
\texttt{\_regenerate\_noise} 在赋值时隐式 cast 回 float64, 物理精度无损.

\begin{lstlisting}[caption={v7.py: lines 2184--2200}]
BATCH_SIZE = 32           # was 64; consumer rate is ~1/slot anyway
NOISE_DTYPE = cp.float32 if GPU_AVAILABLE else np.float32

def __init__(self, buffer, noise_len):
    super().__init__()
    self.buffer = buffer
    self.noise_len = noise_len
    self.stop_event = threading.Event()
    self.daemon = True

def run(self):
    if GPU_AVAILABLE:
        cp.cuda.Device(0).use()
    while not self.stop_event.is_set():
        batch = cp.random.randn(
            self.BATCH_SIZE, 2, self.noise_len).astype(self.NOISE_DTYPE)
        self.buffer.put_batch(batch)
\end{lstlisting}

% ---------------------------------------------------------------------
\bugitem{17}{\texttt{GTBatchSaver} 异步写盘}{\stwo}

\textbf{现象.} \vfour{} 主循环在每个保存点同步调用 \texttt{np.savez\_compressed},
压缩耗时数十至数百毫秒, 期间主循环停滞, 出现 ring buffer 拥塞与 UE drop.

\textbf{\vseven{} 修订.} 后台 daemon thread + \texttt{Queue(maxsize=16)}
异步写盘. queue 满时降级为同步写 (背压机制, 不丢数据);
\texttt{flush\_all()} 先 \texttt{join()} 确保排空.

% ---------------------------------------------------------------------
\bugitem{18}{\texttt{bypass\_copy} 预分配}{\sone}

\textbf{现象.} \vfour{} 每次 bypass copy 调用都执行 \texttt{cp.zeros(...)} 重新分配
GPU 内存, 高频调用下加大 allocator 压力.

\textbf{\vseven{} 修订.} lazy 分配 + 按 \texttt{(nsamps, dst\_nbAnt)} key 缓存
复用 buffer. 同时修复 \vfour{} 中非对称天线场景下
\texttt{dst\_2d[:, min\_ant*2:]} 未初始化的副作用.

% ---------------------------------------------------------------------
\bugitem{19}{UL \texttt{set\_last\_ul\_rx\_ts} 合并}{\sone}

\textbf{现象.} bypass mode 下, \vfour{} 每个 UE 各调用一次
\texttt{ipc\_gnb.set\_last\_ul\_rx\_ts}, 触发一次 \texttt{futex\_wake};
$N$ 个 UE 即 $N$ 次系统调用.

\textbf{\vseven{} 修订.} 收集 \texttt{max\_new\_head}, 循环外仅 wake 一次:
\begin{verbatim}
if _any_bypass:
    self.ipc_gnb.set_last_ul_rx_ts(int(_max_new_head - 1))
\end{verbatim}

% ---------------------------------------------------------------------
\bugitem{20}{\texttt{IPCRingBuffer} 同步移出锁}{\sthree}

\textbf{现象.} \vfour{} 在 \texttt{with s.not\_full:} 锁内调用
\texttt{cp.cuda.Device(0).synchronize()}. CUDA stream 同步耗时取决于
当前 kernel 队列深度 (毫秒级), 期间所有等待该锁的进程被阻塞;
multi-UE 场景下整个 pipeline 被串行化.

\textbf{\vseven{} 修订.} 锁内仅做 index 更新与 \texttt{notify\_all},
\texttt{Device.synchronize()} 移到锁外. 同步完成后 re-notify
防止 consumer 与 producer 间出现竞态.
\begin{verbatim}
with s.not_full:
    # write + index update + notify
    ...
s.not_empty.notify_all()
cp.cuda.Device(0).synchronize()    # outside the lock
with s.not_empty:
    s.not_empty.notify_all()        # re-notify
\end{verbatim}

% =====================================================================
\section{CUDA Graph}
\label{sec:cudagraph}

% ---------------------------------------------------------------------
\bugitem{21}{Capture 模式: Relaxed}{\sone}

\textbf{现象.} \vfour{} 调用 \texttt{stream.begin\_capture()} 默认使用
Global capture mode, 与同时运行的其他 stream (Sionna XLA, NoiseProducer)
存在冲突, capture 偶发失败.

\textbf{\vseven{} 修订.}
\begin{verbatim}
try:
    mode = cp.cuda.runtime.streamCaptureModeRelaxed
except AttributeError:
    mode = 2  # CUDA runtime magic number
stream.begin_capture(mode=mode)
\end{verbatim}

% ---------------------------------------------------------------------
\bugitem{22}{Capture 失败可恢复}{\sone}

\textbf{现象.} \vfour{} 一旦首次 capture 失败, 直接设置
\texttt{self.use\_cuda\_graph = False}, 整个生命周期内不再尝试,
即便后续条件已恢复. 退化为 eager mode 后单 slot 处理时间显著增加.

\textbf{\vseven{} 修订.} 维护 \texttt{\_capture\_failure\_count}
最多重试 5 次; 不直接关闭 graph 路径, 下一 slot 仍会尝试 capture.

% =====================================================================
\section{杂项}
\label{sec:misc}

\begin{center}
\renewcommand{\arraystretch}{1.30}
\begin{tabular}{c p{8.4cm} c}
\toprule
\textbf{编号} & \textbf{修订内容} & \textbf{严重度} \\
\midrule
23 & PEP 604 \texttt{Tuple[...]|None} (Py3.10+) 改为
     \texttt{Optional[Tuple[...]]} (Py3.8 兼容) & \sone \\
24 & CPU 模式下 (\texttt{GPU\_AVAILABLE=False}) \texttt{noise\_dBFS}
     不再被静默丢失, 改用 Python \texttt{float} 计算 \texttt{noise\_std\_abs} & \sone \\
25 & 移除分配后从未使用的 \texttt{self.\_ul\_clip\_3d = cp.zeros(...)} & \sone \\
26 & \texttt{if pl\_linear != 1.0: gpu\_out *= pl\_linear} 改为
     无条件乘 \texttt{cp.float64(pl\_linear)}, 避免 CUDA Graph capture
     时出现 Python 分支 & \sone \\
\bottomrule
\end{tabular}
\end{center}

% =====================================================================
\section{控制面稳定性 (v6 $\to$ v7)}
\label{sec:control}

v6 修复完物理建模后, OAI 在 attach 阶段暴露出新的不稳定现象, 见下表.

\begin{center}
\renewcommand{\arraystretch}{1.25}
\begin{tabular}{l c c}
\toprule
\textbf{配置} & \textbf{结果} & \textbf{结论} \\
\midrule
\texttt{UE\_SPEED=3}, stable 45s   & PARTIAL, srs=0   & 过早进入动态信道, attach/SRS 不稳 \\
\texttt{UE\_SPEED=0}, stable 300s  & OK, srs=1        & 长稳定窗口可完成 SRS 配置 \\
\texttt{UE\_SPEED=3}, stable 300s  & OK, srs=1        & 问题与 \texttt{Speed=3} 本身无关 \\
\texttt{UE\_SPEED=3}, auto trigger & OK, srs=1$\sim$2 & 事件触发可替代固定等待 \\
\bottomrule
\end{tabular}
\end{center}

% ---------------------------------------------------------------------
\bugitem{27}{Attach-Stable Window}{\stwo}

\textbf{现象.} \vsix{} 启动后立即进入动态 \texttt{sample\_times}.
RRC/NAS/SRS 配置阶段未完成时, UL Msg3 收不稳, RRCReconfiguration
无法触发, periodic SRS 因依赖 RRCReconfigurationComplete 而无法启用,
表现为 \texttt{srs=0}.

\textbf{\vsix{} 行为.}
\begin{lstlisting}[caption={v6.py: lines 2119--2120 (启动即动态)}]
sample_times_now = _build_sample_times(batch_idx)
h_c128_norm = generate_fn(sample_times_now)
\end{lstlisting}

\textbf{\vseven{} 修订.} 引入 attach-stable 窗口, 期间将
\texttt{sample\_times} 冻结于 batch 0; 触发后从 0 重新计时, 避免
inter-batch 大相位跳变.
\begin{lstlisting}[caption={v7.py: lines 2186--2196}]
if attach_stable_active:
    # Freeze inter-batch time during attach so RRC/NAS/SRS setup
    # sees a quasi-static channel, then restart dynamic time at 0.
    sample_times_now = _build_sample_times(0)
else:
    if (attach_stable_batches > 0 and
            batch_idx == attach_stable_batches and
            not attach_stable_handoff_logged):
        print("[v7 AttachStable] dynamic channel handoff")
    dyn_batch_idx = batch_idx - dynamic_start_batch
    sample_times_now = _build_sample_times(dyn_batch_idx)
\end{lstlisting}

% ---------------------------------------------------------------------
\bugitem{28}{Auto Handoff}{\sone}

\textbf{背景.} 固定时长稳定窗口对场景敏感: 太短 attach 失败,
太长拉低实验吞吐.

\textbf{\vseven{} 修订.} \texttt{launch\_all\_v7.sh} 后台监控
\texttt{gnb.log}, 检测 \texttt{Received RRCReconfigurationComplete}
后等待 \texttt{ATTACH\_STABLE\_POST\_DELAY=15s},
写入触发文件 \texttt{/tmp/oai\_gpu\_ipc/v7\_dynamic\_enable};
producer 每 10 batch 检查一次该文件存在性.
\begin{lstlisting}[caption={v7.py: lines 2170--2184}]
if attach_stable_active:
    if (attach_stable_mode == "auto" and attach_stable_trigger_file and
            (batch_idx % 10 == 0) and
            _os.path.exists(attach_stable_trigger_file)):
        attach_stable_active = False
        dynamic_start_batch = batch_idx
        attach_stable_handoff_logged = True
        print("[v7 AttachStable] dynamic channel handoff "
              f"triggered by file: {attach_stable_trigger_file}")
    elif batch_idx >= attach_stable_batches:
        attach_stable_active = False
        dynamic_start_batch = batch_idx
        attach_stable_handoff_logged = True
        print("[v7 AttachStable] dynamic channel handoff "
              f"after max stable window ({attach_stable_sec:.1f}s)")
\end{lstlisting}
\texttt{--attach-stable-sec=300} 同时作为 fallback ceiling
(触发文件未出现时强制切换).

% ---------------------------------------------------------------------
\bugitem{29}{P1B 距离推导模式}{\stwo}

\textbf{背景.} \vsix{} 在 P1B 数据可用时强制按 $d \approx \tau_{\min} \cdot c$
推算 BS-UE 距离, 下限 1\,m.

\textbf{现象.} P1B 数据中的 \texttt{tau\_rays} 在 dump 流程中通常以
\textit{相对延迟} 形式存储 (相对于第一径或 frame boundary), 并非
BS-UE 间的绝对传播时延. 实际数据上 $\tau_{\min}$ 经常为亚 ns 级,
$c \cdot \tau_{\min} < 1\,\text{m}$, 全部被 \texttt{tf.maximum(..., 1.0)}
楼至 1\,m, 各 BS-UE 对距离同时塌陷至 1\,m. 该几何下:
\begin{itemize}
  \item 路径损耗值过低, 接收功率超出 OAI gNB 内 AGC 量程;
  \item TA 估算接近 0, 偏离合理蜂窝场景;
  \item RSRP 测量饱和.
\end{itemize}
对应 attach 阶段不稳, \texttt{srs=0}.

\textbf{\vsix{} 行为.}
\begin{lstlisting}[caption={v6.py: lines 2004--2017}]
_default_d = float(cfg.get('default_distance_m', 100.0))
if 'ray_data_stacked' in cfg or 'ray_data' in cfg:
    # Use earliest tap delay across paths to estimate distance
    _tau_first = _tf.reduce_min(tau_rays, axis=-1)
    _tau_first = _tf.squeeze(_tau_first, axis=-1)
    distance_3d = _tau_first * _tf.constant(SPEED_OF_LIGHT, _tau_first.dtype)
    distance_3d = _tf.maximum(distance_3d, 1.0)                 # floor at 1m
else:
    distance_3d = _tf.ones([1, N_BS, N_UE]) * _default_d
\end{lstlisting}

\textbf{\vseven{} 修订.} 引入 \texttt{--p1b-distance-mode \{default, tau\}},
默认 \texttt{default} (使用 \texttt{--default-distance-m=100}m), 仅在显式
请求下走 $\tau$ 推导. 该改动将\textit{几何距离}与\textit{PDP 内部相对延迟}
两个概念解耦, 避免 P1B 数据格式假设传导至 attach 路径损耗.

\begin{lstlisting}[caption={v7.py: lines 2051--2070}]
_default_d = float(cfg.get('default_distance_m', 100.0))
_p1b_distance_mode = str(cfg.get('p1b_distance_mode', 'default')).lower()
_has_p1b = 'ray_data_stacked' in cfg or 'ray_data' in cfg
if _has_p1b and _p1b_distance_mode == 'tau':
    _tau_first = _tf.reduce_min(tau_rays, axis=-1)
    _tau_first = _tf.squeeze(_tau_first, axis=-1)
    distance_3d = _tau_first * _tf.constant(SPEED_OF_LIGHT, _tau_first.dtype)
    distance_3d = _tf.maximum(distance_3d, 1.0)
    print(f"[v7 ChannelProducer] distance_3d derived from P1B tau, "
          f"range=[{float(_tf.reduce_min(distance_3d))}, "
          f"{float(_tf.reduce_max(distance_3d))}] m")
else:
    distance_3d = _tf.ones([1, N_BS, N_UE]) * _default_d
    if _has_p1b:
        print(f"[v7 ChannelProducer] distance_3d=fixed {_default_d} m "
              f"(P1B distance mode={_p1b_distance_mode})")
    else:
        print(f"[v7 ChannelProducer] distance_3d=fallback {_default_d} m "
              f"(no P1B; configure via --default-distance-m)")
\end{lstlisting}

\textbf{物理影响.} 在该 PDP-based 通道模型中, \texttt{distance\_3d}
仅作用于大尺度路径损耗, 不参与小尺度 fading (PDP, Doppler, angular spread)
的计算; 后者完全由 P1B 的 ray data 驱动. 因此固定 100\,m 不引入 PDP 形状失真,
仅恢复合理 path-loss 标度.

% =====================================================================
\section{验证}
\label{sec:verify}

\subsection{5 阶段物理一致性测试 (\texttt{verify\_issue\_2.py})}

\begin{center}
\renewcommand{\arraystretch}{1.25}
\begin{tabular}{c p{6.6cm} l}
\toprule
\textbf{Phase} & \textbf{测试内容} & \textbf{结果} \\
\midrule
A          & 固定 \texttt{sample\_times} 两次调用 producer
           & \texttt{identical=True} (确认 \#2) \\
B (3.5)    & 递增 \texttt{sample\_times}, 保留 \texttt{cf=3.5}
           & \texttt{identical=True} (确认 \#1) \\
B (3.5e9)  & 递增 \texttt{sample\_times}, 修复 \texttt{cf=3.5e9}
           & relative $\Delta=93.5\%$, 信道正确演变 \\
C          & 跨 batch 边界连续性
           & 连续 $2.6\times10^{-7}$ vs 重置 $1.6\times10^{-5}$ ($61.6\times$) \\
D          & \texttt{Speed=0} sanity check
           & \texttt{identical=True}, Doppler 是唯一时变源 \\
\bottomrule
\end{tabular}
\end{center}

\subsection{NMSE / 数据价值初步结论 (\vseven{}, SNR=10\,dB, $UE\_SPEED=3$)}

短窗口数据有效, 长窗口跨越 attach $\to$ 动态切换边界后崩坏.

\begin{center}
\renewcommand{\arraystretch}{1.20}
\begin{tabular}{c c c c}
\toprule
$N$ & NMSE (dB) & RSRP err (dB) & PDP corr \\
\midrule
10  & $-6.35$  & 0.90  & 0.9999 \\
15  & $-7.25$  & 0.75  & 0.9997 \\
20  & $-6.94$  & 0.80  & 0.9994 \\
30  & $-6.83$  & 0.82  & 0.9998 \\
100 & $+14.99$ & 15.13 & 0.2884 \\
\bottomrule
\end{tabular}
\end{center}

建议短窗口取 \texttt{--fixed-n 15} 作为 SNR sweep 起点; 默认 $N=100$
不应作为当前 \vseven{} 动态链路的代表 NMSE.

% =====================================================================
\section{后续待办}
\label{sec:todo}

\begin{itemize}
  \item Auto handoff 后仍可见 \texttt{No SRS signal},
        \texttt{Invalid timing advance offset}, \texttt{UL Failure},
        提示动态切换后链路仍非完全稳定; 计划比较
        \texttt{ATTACH\_STABLE\_POST\_DELAY=15/30},
        \texttt{ATTACH\_STABLE\_TRIGGER=rrc\_reconfig/srs} 的稳定性.
  \item ChannelProducer 仍可见 ring buffer full / dropped, 怀疑影响
        UL 控制面与长窗口 NMSE.
  \item \texttt{IPCRingBuffer wrap-around view vs copy} (v4 issue \#11) 暂未修复;
        当前仅读取场景下风险低.
  \item \texttt{tf.experimental.dlpack} 已废弃, 等 TF 升级后统一迁移.
  \item 全 SNR sweep 前建议先固定 \texttt{--fixed-n 15}, 同时保留 N-axis 曲线
        用于判断每个 SNR 点的有效窗口.
\end{itemize}

\end{document}
